\section*{अध्याय \ref{ch:b3:hydrogeology} --- जलभूविज्ञान}

\begin{solution}{exo:b3:hydrogeology:1}
पहली स्क्रीन पर: $z = 120 - 40 = \qty{80}{m}$, $h = 120 - 12.5 =
\qty{107.5}{m}$, दाब शीर्ष $107.5 - 80 = \qty{27.5}{m}$। दूसरी
स्क्रीन पर: $z = \qty{40}{m}$, $h = \qty{105.0}{m}$, दाब शीर्ष
\qty{65.0}{m}। उथली स्क्रीन पर शीर्ष अधिक है, इसलिए जल नीचे की ओर
बहता है, प्रवणता $(107.5 - 105.0)/40 = \num{0.0625}$ के साथ:
पीज़ोमीटर एक पुनर्भरण क्षेत्र में स्थित हैं। अकेला दाब इसके ठीक
विपरीत संकेत देता।
\end{solution}

\begin{solution}{exo:b3:hydrogeology:2}
$q = K\Delta h/L = 12\times0.30/0.50 = \qty{7.2}{m/d}$। क्षेत्रफल
$\pi\times0.05^2 = \qty{7.85e-3}{m^2}$ है, इसलिए $Q = \qty{0.0565}{m^3/d}$, या
$56.5/24 = \qty{2.36}{L/h}$। \omterm{def:b3:hydrogeology:seepage}{रिसाव वेग} $7.2/0.35 =
\qty{20.6}{m/d}$ है।
\end{solution}

\begin{solution}{exo:b3:hydrogeology:3}
$K = 2.5/86400 = \qty{2.894e-5}{m/s}$, और $k = K\mu/(\rho g) =
\num{2.894e-5}\times\num{1.306e-3}/(999.7\times9.81) =
\qty{3.85e-12}{m^2}$। \qty{20}{\celsius} पर $K$ को
$(998.2/\num{1.002e-3})/(999.7/\num{1.306e-3}) = \num{1.30}$ से गुणा करने पर $K =
\qty{3.25}{m/d}$ मिलता है। यह परिवर्तन लगभग पूरी तरह श्यानता का है।
\end{solution}

\begin{solution}{exo:b3:hydrogeology:4}
$T = 2.5\times180 = \qty{450}{m^2/d}$। जल के प्रसार से
$\num{9e-7}\times180 = \num{1.62e-4}$ मिलता है, जो $S =
\num{2e-4}$ का लगभग \qty{80}{\percent} है; शेष भाग चूना-पत्थर के संपीडन से आता है।
\end{solution}

\begin{solution}{exo:b3:hydrogeology:5}
\cref{prop:b3:hydrogeology:thiem} से, $T = Q\ln(r_2/r_1)/(2\pi(s_1 -
s_2)) = 1500\times\ln10/(2\pi\times1.53) = \qty{359}{m^2/d}$, और $K =
359/40 = \qty{9.0}{m/d}$।
\end{solution}

\begin{solution}{exo:b3:hydrogeology:6}
$Q/4\pi T = 1200/(4\pi\times300) = \qty{0.318}{m}$। 1 दिन बाद, $u =
300^2\times\num{5e-4}/(4\times300\times1) = \num{0.0375}$ और $W = -0.5772
+ 3.2834 + 0.0375 - 0.0004 = \num{2.743}$: $s = \qty{0.873}{m}$। 10
दिन बाद, $u = \num{0.00375}$, $W = -0.5772 + 5.5860 + 0.0038 = \num{5.013}$:
$s = \qty{1.596}{m}$। दस गुना लंबे समय में अवनमन दोगुने से कम
होता है।
\end{solution}

\begin{solution}{exo:b3:hydrogeology:7}
$u \le 0.05$ जब $t \ge r^2S/(4T\times0.05)$। \qty{60}{m} पर: $3600\times
\num{2e-4}/90 = \qty{0.008}{d}$, \qty{11.5}{min}। \qty{500}{m} पर:
\qty{0.556}{d}, \qty{800}{min}, यानी 13 घंटे से अधिक। समय $r^2$ के
अनुपात में बढ़ता है: एक-दो दिन के परीक्षण में सरल रेखा के लिए प्रयुक्त
प्रेक्षण कुएँ पंप किए गए कुएँ से कुछ सौ मीटर के भीतर होने चाहिए।
\end{solution}

\begin{solution}{exo:b3:hydrogeology:8}
मध्य में प्रत्येक कुआँ \qty{200}{m} दूर है: $u = 200^2\times\num{2e-4}/(4
\times450\times10) = \num{4.44e-4}$, $W = \num{7.142}$, $s = 0.382\times
7.142 = \qty{2.728}{m}$, प्रत्येक से, कुल \qty{5.456}{m}। \qty{400}{m} पर
$u = \num{1.78e-3}$, $W = \num{5.757}$, $s =
\qty{2.199}{m}$। प्रत्येक कुआँ दूसरे कुएँ में जल को \qty{2.2}{m} नीचे करता है:
उसके पंप को जल को \qty{2.2}{m} ऊपर उठाना पड़ता है, जिसमें ऊर्जा लगती है,
और पंप के अंतर्ग्रहण के सूखने से पहले बचा अवनमन
कम होता है।
\end{solution}

\begin{solution}{exo:b3:hydrogeology:9}
बाँध के आधार और तल के बीच चार धारारेखाएँ $N_f = 5$
वाहिकाएँ सीमित करती हैं; नौ समविभव रेखाएँ \qty{10}{m} को $N_d = 10$
बूंदों में बाँटती हैं। $q = KHN_f/N_d = 5\times10\times5/10 = \qty{25}{m^3/d}$ प्रति मीटर
बाँध, और \qty{300}{m} के नीचे \qty{7500}{m^3/d}। परिकलित आकृति
गुणक, 0.5 के बजाय 0.512, प्रति मीटर \qty{25.6}{m^3/d} देता।
\end{solution}

\begin{solution}{exo:b3:hydrogeology:10}
$k = na^2/8 = 0.30\times(\num{5e-5})^2/8 = \qty{9.4e-11}{m^2}$, और $K =
k\rho g/\mu = \num{9.375e-11}\times998.2\times9.81/\num{1.002e-3} =
\qty{9.2e-4}{m/s}$, \qty{79}{m/d}, जो स्वच्छ बालू की परास, 0.3 से
\qty{150}{m/d} के भीतर है। प्रतिरूप में पथों की घुमावदारी छूटी है,
जो सीधी नलिकाओं से लंबे होते हैं (इससे $K$ घटता है), रंध्रों के बीच के कंठ,
जो रंध्रों से संकरे होते हैं (निम्न), वे रंध्र जो अंधे समाप्त
होते हैं और प्रवाह नहीं ले जाते (निम्न), और रंध्र-आकारों का
विस्तार: चूँकि नलिका प्रवाह को $a^4$ के अनुपात में ले जाती है,
सबसे चौड़े पथ अधिकांश जल ले जाते हैं, और विस्तार $K$ को माध्य आकार की नलिकाओं वाले मान से ऊपर ले जाता है।
\end{solution}

\begin{solution}{exo:b3:hydrogeology:11}
(क)~लघुगणक शून्य हो जाता है जब $2.25\,Tt = r^2S$, अर्थात् $r_0 =
\sqrt{2.25\,Tt/S}$ पर। (ख)~परिरुद्ध और अपरिरुद्ध जलभृतों के लिए:
\[
r_0 = \sqrt{\frac{2.25\times450\times30}{\num{2e-4}}} = \qty{12.3}{km},
\qquad r_0 = \sqrt{\frac{2.25\times450\times30}{0.2}} = \qty{390}{m},
\]
$\sqrt{1000} = 31.6$ गुना छोटा। (ग)~दाब विसरण $T/S$ के साथ फैलता है,
जो $S$ छोटा होने पर बहुत बड़ा होता है: परीक्षण के भीतर शंकु
भ्रंशों, जलभृत के अनावरणों और नदियों तक पहुँचता है, और अवनमन सरल
रेखा से हट जाता है, जो भूजलविज्ञानी को बताता है
कि किसी सीमा से सामना हो चुका है।
\end{solution}

\begin{solution}{exo:b3:hydrogeology:12}
(क)~जल के साथ चलने वाले फ़्रेम से देखने पर अग्र भाग का मानक विचलन
$\sigma_x = \sqrt{2Dt} = \sqrt{2\alpha_Lvt}$ के साथ फैला होता है। यह
$x$ को $t = x/v$ के आसपास पार करता है, जब $\sigma_x = \sqrt{2\alpha_Lx}$; वेग
$v$ पर दूरी में $\sigma_x$ का फैलाव समय में $\sigma_t = \sigma_x/v$ का होता है।
(ख)~\qty{16}{\percent} और \qty{84}{\percent} समय
माध्य के दोनों ओर एक-एक मानक विचलन पर पड़ते हैं: $\sigma_t = (247.2 -
152.6)/2 = \num{47.3}\text{ दिन}$। तब $\alpha_L = v^2\sigma_t^2/(2x) =
0.25\times47.3^2/200 = \qty{2.8}{m}$। (वक्र \qty{3}{m} से परिकलित किया गया था:
प्रसामान्य सन्निकटन इसे थोड़ा कम आँकता है।) (ग)~$R =
1 + 1.9\times0.32/0.30 = \num{3.03}$: विलेय जल की गति के एक-तिहाई
वेग से चलता है।
\end{solution}

\begin{solution}{pb:b3:hydrogeology:1}
\textbf{1.} शेल $363.2 - 282.5 = \qty{80.7}{m}$; चूना-पत्थर $282.5 -
102.5 = \qty{180}{m}$।

\textbf{2.} $h = 363.2 - 45.0 = \qty{318.2}{m}$; ऊपरी सतह पर दाब शीर्ष
$318.2 - 282.5 = \qty{35.7}{m}$। जल चूना-पत्थर की ऊपरी सतह से
\qty{35.7}{m} ऊपर उठता है, जिसे शेल नीचे दबाए रखता है: जलभृत परिरुद्ध है।

\textbf{3.} $h_2 - h_1 = 550a - 100b = -1.62$ और $h_3 - h_1 = 250a
- 600b = -1.44$ से: $a = \num{-2.715e-3}$ और $b = \num{1.269e-3}$ प्रति
मीटर।

\textbf{4.} $|\nabla h| = \sqrt{a^2 + b^2} = \num{0.0030}$। जल
$-\nabla h = (\num{2.715e-3}, \num{-1.269e-3})$ के अनुदिश बहता है, पूर्व और
दक्षिण की ओर: दिगंश $\ang{180} - \arctan(2.715/1.269) = \ang{115}$।

\textbf{5.} लगभग \qty{10}{min} से पहले के बिंदु रेखा के ऊपर पड़ते हैं: वहाँ
$u > 0.05$ है और थीस वक्र अपने लघुगणकीय सन्निकटन से ऊपर है
($W > -\gamma - \ln u$, जैसा कुएँ के फलन के आरेख से दिखता है)।

\textbf{6.} \qty{15}{min} के बाद के बिंदुओं पर आधारित सर्वोत्तम रेखा
\qty{100}{min} पर \qty{1.75}{m} और \qty{1000}{min} पर \qty{2.63}{m} से
होकर गुजरती है: प्रति लघुगणक चक्र $\Delta s = \qty{0.87}{m}$।

\textbf{7.} $T = 2.3\times2160/(4\pi\times0.872) = \qty{453}{m^2/d}$।

\textbf{8.} रेखा \qty{100}{min} पर \qty{1.75}{m} से प्रति चक्र
\qty{0.87}{m} गिरती है: शून्य $1.75/0.872 = 2.01$ चक्र पहले पहुँचती है,
यानी $t_0 = \qty{0.98}{min} = \qty{6.83e-4}{d}$।

\textbf{9.} $S = 2.25\,Tt_0/r^2 = 2.25\times453\times\num{6.83e-4}/3600 =
\num{1.9e-4}$, परिरुद्ध जलभृतों की परास के भीतर।

\textbf{10.} \qty{15}{min} पर, $u = 3600\times\num{1.9e-4}/(4\times453
\times0.0104) = \num{0.036}$, और उसके बाद इससे कम।

\textbf{11.} $K = 453/180 = \qty{2.5}{m/d}$, एक दरारयुक्त चूना-पत्थर।

\textbf{12.} $u = 3600\times\num{2e-4}/(4\times450\times2) =
\num{2.0e-4}$, $W = -0.5772 + 8.5172 + 0.0002 = \num{7.940}$, $s =
0.382\times7.940 = \qty{3.03}{m}$, जैसा अभिलेखित है।

\textbf{13.} $u = 10^6\times\num{2e-4}/(4\times450\times30) =
\num{0.0037}$, $W = \num{5.02}$, $s = \qty{1.92}{m}$।

\textbf{14.} $r_0 = \sqrt{2.25\times450\times30/(\num{2e-4})} =
\qty{12.3}{km}$, जो क्षेत्र की चौड़ाई से अधिक है। शंकु चूना-पत्थर के
किनारों, उसके अनावरणों और भ्रंश तक पहुँच चुका है, जो उसे कम पारगम्य
चट्टान के विरुद्ध रोक सकता है: थीस का अनंत जलभृत अब इसका वर्णन नहीं
करता, और केवल अवनमन के अभिलेख से पता चलेगा कि कोई सीमा कुएँ को जल
देती है या प्रवाह रोकती है।

\textbf{15.} $2160\times30 = \qty{64800}{m^3}$। यह संचयन से आता है:
जल के प्रसार और चूना-पत्थर के संपीडन से मुक्त हुआ जल,
शंकु के पूरे क्षेत्र पर। शंकु पर $S\,s$ का समाकलन
$Qt$ वापस देता है।

\textbf{16.} $u = 10^6\times0.18/(4\times450\times30) = \num{3.33}$, $W
= \num{0.0086}$, $s = \qty{3}{mm}$: फ़ार्म को इसका पता नहीं चलेगा। पंपिंग
कुएँ के पास के रंध्रों से जल खींचेगी, न कि
पूरे जलभृत के दाब से।

\textbf{17.} $q = 2.5\times0.0030 = \qty{0.0075}{m/d}$; $v = 0.0075/0.18
= \qty{0.042}{m/d}$, प्रति वर्ष \qty{15}{m}।

\textbf{18.} $200/0.0417 = \num{4800}\text{ दिन}$, \num{13.1}~वर्ष; प्रदूषक,
\num{32.9}~वर्ष।

\textbf{19.} त्रिज्या $r$ के बेलन से होकर प्रवाह $Q$ क्षेत्रफल $2\pi rb$ को
पार करता है, जिसका रंध्र-भाग $n_e$ है: $\mathrm dr/\mathrm
dt = -Q/(2\pi rbn_e)$, और $2\pi bn_er\,\mathrm dr = -Q\,
\mathrm dt$ का $r$ से 0 तक समाकलन करने पर $t = \pi n_ebr^2/Q$ मिलता है। यहाँ $\pi\times0.18
\times180\times200^2/2160 = \num{1885}\text{ दिन}$, जल के लिए \num{5.2}\text{ वर्ष}
और प्रदूषक के लिए $2.5\times5.2 = \num{12.9}\text{ वर्ष}$: कुआँ
उसे दोगुने से भी अधिक जल्दी पहुँचा देता है।

\textbf{20.} $\alpha_L = \qty{20}{m}$ और $\sqrt{2\times20\times200} =
\qty{89}{m}$। सांद्रता अपने अंतिम मान के \qty{16}{\percent} तक
माध्य अग्रभाग से \qty{89}{m} आगे पहुँच जाती है: यह स्तर कुएँ तक पहुँचता है
तब, जब माध्य अग्रभाग अभी \qty{89}{m} दूर है, $(111/200)^2\times12.9 = \num{4.0}\text{ वर्ष}$ पर,
माध्य समय का एक-तिहाई। पेयजल के लिए पहले के निशान मायने रखते हैं,
माध्य नहीं।

\textbf{21.} पक्ष में: कुआँ प्रदूषक को पकड़ लेता है, इसलिए वह और नीचे की
प्रवणता की ओर नहीं फैल सकता, और जल को जहाँ से पंप किया जाता है वहीं उपचारित किया जा सकता है।
विपक्ष में: पंपिंग प्रदूषक को किसी आपूर्ति कुएँ तक उससे पहले खींच लाती है,
जितना वह अन्यथा कहीं भी पहुँचता, और जल को दशकों तक निगरानी में रखना
और उपचारित करना पड़ता है।

\textbf{22.} वार्डन चूना-पत्थर की \omterm{def:b3:hydrogeology:storage}{पारगमनता} $T \approx
\textbf{\qty[text-series-to-math=true]{450}{m^2/d}}$ है (परीक्षण से
\qty{453}{m^2/d})।
\end{solution}
